\begin{proposition}[A-L07: Proved; contractible horizontal parameters still need a strict lift]
\label{AL07}
There is an admissible integral theory $A$ with finite free
terms and support $q\leq w$, a degreewise free contractible
simplicial module $M$, and a simplicial linear map
$\alpha:M\to A^{2,1}$ for which the cycle map
\[
 I_N\otimes_{\mathbb Z}\Gamma^2_{\mathbb Z}(M)
       \longrightarrow Z^4(A),\qquad
 a\otimes\gamma_2(m)\longmapsto a\,\gamma_2(d\alpha(m))
 \tag{AL.36}
\]
has no weight-preserving, simplicial $S$-linear primitive
in $A^3$, for every $N\geq1$.
Here $\Gamma^2(M)$ is placed in weight four and
$I_N$ has its coefficient weights.
The target resolves $S$ in all form degrees in every
weight $w\geq5$, uniformly in the simplex.
\end{proposition}

\begin{proof}
Use the PD differential algebra $B$ of
Theorem~\ref{AL03}, including its elements $x,y,u,\vartheta$,
the map $\pi':F'\to S\vartheta$, and its augmentation
$B\to S$.
Replace the vertex-function simplicial algebra $\Phi$
there by the following explicitly defined simplicial
algebra:
\[
 C_n=\mathbb Z^{\{(i,j):0\leq i\leq j\leq n\}}.
 \tag{AL.37}
\]
Multiplication is coordinatewise.  For an order-preserving
map $\eta:[m]\to[n]$, the map $C_n\to C_m$ sends a
function $f$ to the function
$(i,j)\mapsto f(\eta(i),\eta(j))$.
Composition of these functions verifies the simplicial
identities and preservation of all ring operations.
Faces restrict along the omitted vertex; degeneracies
repeat the corresponding coordinates.

We compute the entire normalized complex with the
convention $N_nC=\bigcap_{i=1}^n\ker d_i$.
Let $\epsilon_{ij}$ be the indicator function of the
coordinate $(i,j)$.  Then
\[
 N_0C=\mathbb Z\epsilon_{00},\quad
 N_1C=\mathbb Z\epsilon_{01}\oplus\mathbb Z\epsilon_{11},
 \quad N_2C=\mathbb Z\epsilon_{12},\quad
 N_nC=0\quad(n\geq3).
 \tag{AL.38}
\]
Indeed membership in $\ker d_i$ requires a function to
vanish on every pair not containing the vertex $i$.
A surviving pair must therefore contain every vertex
$1,\ldots,n$.  For $n=1,2$ this gives exactly the listed
pairs, and for $n\geq3$ no pair has that property.
The normalized differential $d_0$ sends
$\epsilon_{12}$ to $\epsilon_{01}$,
$\epsilon_{11}$ to $\epsilon_{00}$, and
$\epsilon_{01}$ to zero.  Thus $NC$ is the direct sum
of the identity disks in homological degrees $(1,0)$
and $(2,1)$.  Explicitly the homotopy sends
$\epsilon_{00}$ to $\epsilon_{11}$ and
$\epsilon_{01}$ to $\epsilon_{12}$ and is zero on the
other listed basis elements.  Its differential
commutator is the identity.  Dold--Kan transports
this contraction to an additive simplicial contraction.

Put $R_n=G_n\otimes_SB$ as in Theorem~\ref{AL03}, and set
\[
 A_n=\mathop{\prod\nolimits_{G_n}}_{0\leq i\leq j\leq n}R_n
      =G_n\oplus(G_n\otimes_SB_+)\otimes_{\mathbb Z}C_n .
 \tag{AL.39}
\]
All factors have the common augmentation to $G_n$.
The first description defines products and divided
powers coordinatewise.  The simplicial maps apply
the corresponding map of $G$ and the coordinate
rule \textup{(AL.37)}.
Consequently this is an actual simplicial strict PD
differential algebra over $S$, with the specified
coefficient map and augmentation.
The proof of Theorem~\ref{AL03} applies with this
explicit contraction of $C$: every added form-weight
row is contractible, and the vertical $S$-linear
contraction $G_n\simeq S$ computes its sole additional
cohomology as
$(\mathbb Z/2)^{\binom{n+2}{2}}$ in form degree four,
weight four.  This proves the stated uniform
admissibility.  The finite coordinate count, the
finite free pieces of $B,G_n$, and their established
support prove the finite freeness and $q\leq w$.

Let $M=\operatorname{DK}(D^2(\mathbb Z))$, where
$D^2(\mathbb Z)$ is the normalized identity disk
$\mathbb Z\to\mathbb Z$ in homological degrees
two and one.  The chain map sending its top
generator $m_2$ to $\epsilon_{12}$ and its bottom
generator to $\epsilon_{01}$ defines a simplicial
linear map $f:M\to C$ by the exact Dold--Kan
equivalence.  Each $M_n$ is free, and $M$ is
contractible by the identity-disk contraction.
Define
\[
 \alpha_n(m)=(0,\,1\otimes x\otimes f_n(m))
                  \in A_n^{2,1}.
 \tag{AL.40}
\]
This is simplicial and linear.  Its differential
has $y$ in place of $x$, and its divided square
has $\vartheta$ times the coordinatewise square
of $f_n(m)$.  The PD universal property makes
this quadratic expression a linear map from
$\Gamma^2(M)$; it is closed since $dy=0$.
This defines \textup{(AL.36)} in its stated category.

Suppose a primitive of \textup{(AL.36)} existed.
In simplex degree two, evaluate it on
$a\otimes\gamma_2(m_2)$.
Project its output to the coordinate $(1,2)$ in
\textup{(AL.39)}, and then apply
$G_2\otimes_SB\to G_0\otimes_SB=B$ induced by
any chosen vertex map $G_2\to G_0=S$.
These are weight-preserving $S$-linear chain maps,
and the coordinate value of $f_2(m_2)$ is one.
Their composite would therefore give an $S$-linear
map $I_N\langle4\rangle\to B^3=F'$ whose differential
is $a\mapsto a\vartheta$.
The reduction-modulo-two calculation at the end of
Theorem~\ref{AL03} rules out exactly this map for
every $N$.  This contradiction proves the claim.
\end{proof}

The obstruction in this proposition is to a specified
map and a specified primitive, although its parameter
is already contractible.  The coefficient-scaled
maps of Theorem~\ref{AL05} satisfy their additional
explicit formula and are unaffected by this
counterexample.  The section $G\to A$ and
augmentation $A\to G$ in \textup{(AL.39)} again
have composite the identity.
